\section{Conclusion}\label{sec:conclusion}

Preserving an external benchmark is evidence about the link between benchmark and method. For synthetic cultural agents, which compile declared aggregate preference signs into choice policies that read survey items without naming a group, we tested each link of the chain from anchor to instrument to respondents separately.

Between anchor and instrument the method holds in a limited sense. Training installs the synthetic labels on pairs the adapters never saw, and on trust the adapters keep the anchor's order of countries across a change of scenario bank. A country name does not keep that order, and a prompt stating the same signs does, so the declared signs carry the trust order and the evidence cannot credit the trained weights over a prompt. The adapters do not keep the declared dimensions separate. Between instrument and respondents the evidence is weak. A preregistered screen of held-out items, read against a retrospective criterion, gives a small signal that the registered comparisons do not separate from careful readers of the questionnaire, no dimension passes both the second-bank check and the screen, and the adapters' order of countries does not detectably agree with the respondents' composite. Keeping the anchor's order of countries and approaching respondents' answer distributions are different questions, and on a descriptive reading the adapters keep the first and lie farther from the survey's answer distributions than the base model, whereas a country-named prompt lies closer with no trust order.

The method is an inspectable, group-level reading of candidate items that a researcher may use beside human pretesting, as a complement and never in its place. The evaluation protocol applies to other conditioning methods and other anchors. The test that would license more is prospective: probe-selected items fielded to people after the selection and a human-only comparison rule have both been declared. Until that test is run, pre-piloting remains untested.
